\section{Kramers--Wannier duality squared}%
\label{sec:asymex_kw}

Consider $N\geq1$ qubits on a periodic chain, with sites
$I_N=\{0,\ldots,N-1\}$ and all site indices understood modulo $N$.
The Hilbert space $\mathcal H_N=(\C^2)^{\otimes N}$ has computational
basis $\{\ket{b}:b\in\{0,1\}^{I_N}\}$. Write $\bar b$ for the
componentwise complement of $b$, and $b\oplus e_j$ for the configuration
obtained by flipping its bit at site $j$. The Pauli operators act by
\begin{align}
	X_j\ket{b} & =\ket{b\oplus e_j},
	\notag                           \\
	Z_j\ket{b} & =(-1)^{b_j}\ket{b}.
	\label{eq:asymex_kw_pauli}
\end{align}
They are self-adjoint involutions, satisfy $X_jZ_j=-Z_jX_j$, and commute
at distinct sites. Define global spin flip $\eta$ and one-site left
translation $T$ by
\begin{align}
	\eta\ket{b} & =\ket{\bar b},
	\notag                       \\
	T\ket{b}    & =\ket{\tau b},
	\qquad (\tau b)_j=b_{j+1}.
	\label{eq:asymex_kw_translation}
\end{align}
Thus $\eta=\prod_jX_j$, $\eta^2=\Id$, $T^N=\Id$,
$T^\dagger=T^{-1}$, and $T\eta=\eta T$. In this convention,
$TX_jT^{-1}=X_{j-1}$ and $TZ_jT^{-1}=Z_{j-1}$.

For real couplings $J,h$, the periodic transverse-field Ising Hamiltonian is
\begin{align}
	H_N(J,h) & =-J\sum_{j\in I_N}Z_jZ_{j+1}-h\sum_{j\in I_N}X_j.
	\label{eq:asymex_kw_ising}
\end{align}
The sum includes one bond per site. In particular, the bond at $N=1$
is $Z_0^2=\Id$, and at $N=2$ the two bond terms both equal $Z_0Z_1$.
Kramers--Wannier duality exchanges these bond and field terms. On a
periodic spin chain it is a noninvertible operator, rather than a
unitary change of variables on the full Hilbert space
\cite{Aasen2016Topological, Seiberg2024Majorana}. Its matrix product
representation gives a concrete example of the noninvertible MPO
symmetries discussed in Ref.~\cite{Cirac2021Matrix}.

\begin{definition}[The duality kernel and the two translation tensors]
	\label{def:asymex_kw}
	\lean{KWExample.kwTensor, KWExample.shiftTensor, KWExample.flipShiftTensor,
		KWExample.kwSquare, KWExample.kwSquareTargets}
	\leanok
	\uses{def:mpo_tensor, def:mpo_to_mps, def:mpdo_operator_product}
	With the outgoing label $s'$ written first, the duality tensor is
	\begin{align}
		W^{00} & =\begin{pmatrix}1&1\\0&0\end{pmatrix},
		       &
		W^{01} & =\begin{pmatrix}0&0\\1&1\end{pmatrix},
		\notag                                           \\
		W^{10} & =\begin{pmatrix}1&-1\\0&0\end{pmatrix},
		       &
		W^{11} & =\begin{pmatrix}0&0\\-1&1\end{pmatrix}.
		\notag
	\end{align}
	The translation tensor is $T^{s's}=E_{s,s'}$ and the flipped
	translation tensor is $(\eta T)^{s's}=E_{s,1-s'}$, where $E_{pq}$ is
	the matrix unit with a single entry $1$ in row $p$ and column $q$.
	The same symbols $T$ and $\eta T$ denote their periodic operators.
	The stacked product tensor is
	$B^{s's}=\sum_{m=0}^1W^{s'm}\otimes W^{ms}$, of bond dimension four.
	Its two target tensors are $2T$ and $2\eta T$.
\end{definition}

Define the raw periodic operator $K_N$ by the ordinary virtual trace,
without dividing by the bond dimension:
\begin{align}
	(K_N)_{ab}
	 & =\tr(W^{a_0b_0}\cdots W^{a_{N-1}b_{N-1}}).
	\label{eq:asymex_kw_periodic}
\end{align}
We also write $D=K_N$ for this raw operator, retaining $D$ in the
square identity below. The normalized operators will be denoted
$\mathcal D_N$ and $U_N$, not $D$.

\begin{theorem}[Periodic kernel]
	\label{lem:asymex_kw_kernel}
	\lean{KWExample.kwTensor_apply,
		KWExample.kwTensor_mpo_eq_prod,
		KWExample.kwTensor_mpo_eq_pow_sum}
	\leanok

	For all configurations $s',s$,
	\begin{align}
		\bra{s'}D\ket{s}
		 & =\prod_{j\in I_N}(-1)^{s'_j(s_j+s_{j+1})}.
		\label{eq:asymex_kw_kernel}
	\end{align}
\end{theorem}

\begin{proof}
	\leanok
	The local entries are $(W^{ab})_{uv}=\delta_{u,b}(-1)^{a(b+v)}$.
	Expanding the trace over closed virtual configurations $g$ gives
	\begin{align}
		(K_N)_{ab}
		 & =\sum_{g\in\{0,1\}^{I_N}}
		\prod_{j\in I_N}\delta_{g_j,b_j}
		                (-1)^{a_j(b_j+g_{j+1})}
		=\prod_{j\in I_N}(-1)^{a_j(b_j+b_{j+1})}.
		\notag
	\end{align}
	Only the configuration $g=b$ contributes.
\end{proof}

For subsequent calculations, put $q(a,b)=\sum_j a_j(b_j+b_{j+1})$,
so $(K_N)_{ab}=(-1)^{q(a,b)}$. Exponents may be calculated modulo two.
Unless indicated otherwise, site sums and products run over $I_N$.

\begin{theorem}[Spin-flip absorption and noninvertibility]
	\label{thm:asymex_kw_absorption}
	\lean{KWExample.kwTensor_mpo_flip_input,
		KWExample.kwTensor_mpo_flip_output,
		KWExample.kwTensor_mpo_mul_spinFlip,
		KWExample.spinFlip_mul_kwTensor_mpo,
		KWExample.kwTensor_mpo_mulVec_eq_zero_of_odd,
		KWExample.kwTensor_mpo_not_isUnit}
	\leanok
	\uses{lem:asymex_kw_kernel}
	For every $N>0$, the raw operator satisfies
	$K_N\eta=\eta K_N=K_N$. If $\eta v=-v$, then $K_Nv=0$.
	In particular, $K_N$ is not invertible.
\end{theorem}
\begin{proof}
	\leanok
	Complementing either input or output leaves the kernel unchanged:
	$q(a,\bar b)\equiv q(\bar a,b)\equiv q(a,b)\pmod2$.
	This proves both absorption identities. For an odd vector,
	$K_Nv=K_N\eta v=-K_Nv$, so $K_Nv=0$.
	The nonzero vector $\ket{0}^{\otimes N}-\ket{1}^{\otimes N}$
	has odd spin flip and witnesses noninvertibility.
\end{proof}

\begin{lemma}[Local duality and coupling exchange]
	\label{lem:asymex_kw_local}
	For every site $j$,
	\begin{align}
		K_NX_j        & =Z_{j-1}Z_jK_N,
		\notag                          \\
		K_NZ_jZ_{j+1} & =X_jK_N.
		\label{eq:asymex_kw_local}
	\end{align}
	Consequently, the Hamiltonians with exchanged couplings satisfy
	\begin{align}
		K_NH_N(J,h) & =H_N(h,J)K_N.
		\label{eq:asymex_kw_exchange}
	\end{align}
	In particular, $K_N$ commutes with $H_N(J,J)$.
\end{lemma}

\begin{proof}
	Flipping one input or output bit gives
	\begin{align}
		q(a,b\oplus e_j) & \equiv q(a,b)+a_{j-1}+a_j\pmod2,
		\notag                                              \\
		q(a\oplus e_j,b) & \equiv q(a,b)+b_j+b_{j+1}\pmod2.
		\notag
	\end{align}
	These congruences prove~(\ref{eq:asymex_kw_local}) entrywise.
	They remain valid when the two indicated sites coincide.
	Summing the local relations and reindexing the bonds yields
	\begin{align}
		K_NH_N(J,h)
		 & =\left(-J\sum_jX_j-h\sum_jZ_{j-1}Z_j\right)K_N
		=H_N(h,J)K_N.
		\notag
	\end{align}
\end{proof}

The local identities also exchange strings and endpoint operators.
For consecutive sites $j_r=j_0+r$ and $1\leq m\leq N$,
\begin{align}
	K_N\prod_{r=0}^{m-1}X_{j_r}
	 & =Z_{j_0-1}Z_{j_{m-1}}K_N,
	\notag                                        \\
	K_NZ_{j_0}Z_{j_m}
	 & =\left(\prod_{r=0}^{m-1}X_{j_r}\right)K_N.
	\label{eq:asymex_kw_strings}
\end{align}
Indeed, iterating~(\ref{eq:asymex_kw_local}) cancels every interior
factor $Z_j^2$. At $m=N$ the endpoints coincide, and these relations
give $K_N\eta=\eta K_N=K_N$. They are intertwining identities; no
inverse of $K_N$ is involved.

\begin{lemma}[Adjoint and Gram matrices]
	\label{lem:asymex_kw_gram}
	The raw operator satisfies
	\begin{align}
		TK_N        & =K_NT,
		\notag                    \\
		K_N^\dagger & =T^{-1}K_N,
		\label{eq:asymex_kw_adjoint}
	\end{align}
	and
	\begin{align}
		K_N^\dagger K_N=K_NK_N^\dagger & =2^N(\Id+\eta).
		\label{eq:asymex_kw_gram}
	\end{align}
\end{lemma}

\begin{proof}
	Cyclic reindexing gives $q(\tau a,\tau b)=q(a,b)$ and
	$q(b,a)=q(\tau a,b)$. Since the kernel is real, these identities
	imply~(\ref{eq:asymex_kw_adjoint}).

	For a binary configuration $c$, the character sum factors as
	\begin{align}
		\sum_a(-1)^{\sum_ja_jc_j}
		 & =\prod_j(1+(-1)^{c_j})
		=\begin{cases}
			 2^N, & c=\mathbf0,    \\
			 0,   & c\neq\mathbf0.
		 \end{cases}
		\label{eq:asymex_kw_character}
	\end{align}
	Here $\mathbf0$ and $\mathbf1$ denote the two constant configurations.
	For input configurations $b,c$, put $d=b\oplus c$. Then
	\begin{align}
		(K_N^\dagger K_N)_{bc}
		 & =\sum_a(-1)^{\sum_ja_j(d_j+d_{j+1})}
		=2^N(\delta_{b,c}+\delta_{b,\bar c}).
		\notag
	\end{align}
	Equation~(\ref{eq:asymex_kw_character}) imposes $d_j=d_{j+1}$ at
	every site. Such a configuration is constant on the cycle, and
	the two constant configurations are distinct because $N\geq1$.
	For the other Gram matrix, put $d=a\oplus c$ and sum over inputs:
	\begin{align}
		(K_NK_N^\dagger)_{ac}
		 & =\sum_b(-1)^{\sum_jb_j(d_j+d_{j-1})}
		=2^N(\delta_{a,c}+\delta_{a,\bar c}).
		\notag
	\end{align}
	This proves~(\ref{eq:asymex_kw_gram}).
\end{proof}

Define the spin-flip sectors $\mathcal H_N^\pm=\{v:\eta v=\pm v\}$ and
their orthogonal projections $P_\pm=(\Id\pm\eta)/2$. Two useful
normalizations of the duality operator are
\begin{align}
	\mathcal D_N & =2^{-N/2}K_N,
	\notag                                               \\
	U_N          & =2^{-(N+1)/2}K_N=\mathcal D_N/\sqrt2.
	\label{eq:asymex_kw_normalizations}
\end{align}

\begin{lemma}[Kernel, image, and partial isometry]
	\label{lem:asymex_kw_sectors}
	The kernel and image of the raw duality operator are
	\begin{align}
		\ker K_N  & =\mathcal H_N^-,
		\notag                       \\
		\ran K_N  & =\mathcal H_N^+,
		\notag                       \\
		\rank K_N & =2^{N-1}.
		\label{eq:asymex_kw_sectors}
	\end{align}
	The normalized operators satisfy
	\begin{align}
		\mathcal D_N^\dagger\mathcal D_N
		=\mathcal D_N\mathcal D_N^\dagger & =2P_+,
		\notag                                     \\
		U_N^\dagger U_N=U_NU_N^\dagger    & =P_+.
		\label{eq:asymex_kw_partial}
	\end{align}
	Thus $U_N$ is a partial isometry with initial and final projection
	$P_+$. It vanishes on $\mathcal H_N^-$ and restricts to a unitary
	operator on $\mathcal H_N^+$.
\end{lemma}

\begin{proof}
	Equation~(\ref{eq:asymex_kw_gram}) gives
	$\|K_Nv\|^2=2^{N+1}\|P_+v\|^2$, which identifies the kernel.
	The absorption relation $\eta K_N=K_N$ puts the image in
	$\mathcal H_N^+$. Conversely, for $y=P_+y$,
	$K_N(2^{-(N+1)}K_N^\dagger y)=y$, so every vector in that sector
	is in the image. For configurations with $b_0=0$, the vectors
	$(\ket{b}\pm\ket{\bar b})/\sqrt2$ form orthonormal bases of the
	two sectors. Each sector therefore has dimension $2^{N-1}$.
	Rescaling~(\ref{eq:asymex_kw_gram}) proves
	(\ref{eq:asymex_kw_partial}).
\end{proof}

Each $H_N(J,h)$ preserves both spin-flip sectors: the field terms
commute with $\eta$, and the two signs acquired by a bond term cancel.
Thus $U_N$ implements a unitary equivalence between $H_N(J,h)$ and
$H_N(h,J)$ on $\mathcal H_N^+$, but not on the full Hilbert space.

The product state and the Greenberger--Horne--Zeilinger (GHZ) state
give elementary finite-chain examples of the action. Use the
unnormalized vectors
\begin{align}
	p_N & =(\ket{0}+\ket{1})^{\otimes N}=\sum_b\ket{b},
	\notag                                              \\
	g_N & =\ket{\mathbf0}+\ket{\mathbf1}.
	\label{eq:asymex_kw_states}
\end{align}
For $h>0$, $p_N$ is proportional to the unique ground state of
$H_N(0,h)$. For $J>0$ and $N\geq2$, the ground space of $H_N(J,0)$ is
$\spn_{\C}\{\ket{\mathbf0},\ket{\mathbf1}\}$, and $g_N$ spans its
spin-flip-even subspace. These facts follow respectively by minimizing
each field term $-hX_j$ and each diagonal bond term $-JZ_jZ_{j+1}$.
They do not require a thermodynamic limit.

\begin{lemma}[Product-state and GHZ actions]
	\label{lem:asymex_kw_states}
	For every $N\geq1$,
	\begin{align}
		K_Np_N                              & =2^Ng_N,
		\notag                                         \\
		K_N\ket{\mathbf0}=K_N\ket{\mathbf1} & =p_N,
		\notag                                         \\
		K_Ng_N                              & =2p_N.
		\label{eq:asymex_kw_state_actions}
	\end{align}
	Writing $\widehat p_N=2^{-N/2}p_N$ and
	$\widehat g_N=2^{-1/2}g_N$, one has
	$U_N\widehat p_N=\widehat g_N$ and
	$U_N\widehat g_N=\widehat p_N$.
\end{lemma}

\begin{proof}
	Reindexing the exponent and using~(\ref{eq:asymex_kw_character})
	gives
	\begin{align}
		(K_Np_N)_a
		 & =\sum_b(-1)^{\sum_jb_j(a_j+a_{j-1})}
		=2^N(\delta_{a,\mathbf0}+\delta_{a,\mathbf1}).
		\notag
	\end{align}
	Both constant inputs have phase one against every output
	configuration, so each maps to $p_N$. Finally,
	$\|p_N\|^2=2^N$ and $\|g_N\|^2=2$ give the normalized statements.
\end{proof}

In particular, the antisymmetric combination
$\ket{\mathbf0}-\ket{\mathbf1}$ is annihilated. These finite-chain
identities alone assert neither spontaneous symmetry breaking nor a
thermodynamic spectral-gap property.

The square follows already from the adjoint and Gram identities.
With $D=K_N$ as above,
\begin{align}
	D^2 & =2^N(\Id+\eta)T.
	\label{eq:asymex_kw_fusion}
\end{align}
Indeed, $T^{-1}K_N^2=K_N^\dagger K_N$ by
(\ref{eq:asymex_kw_adjoint}), and
(\ref{eq:asymex_kw_gram}) gives the result. The corresponding normalized
identities are
\begin{align}
	\mathcal D_N^2 & =(\Id+\eta)T,
	\notag                         \\
	U_N^2          & =P_+T,
	\notag                         \\
	U_N^{2N}       & =P_+.
	\label{eq:asymex_kw_normalized_square}
\end{align}
The last identity uses $P_+T=TP_+$ and $T^N=\Id$.
The two summands $T$ and $\eta T$ act on the same Hilbert space.
On $\mathcal H_N^+$ they add, giving $K_N^2=2^{N+1}T$; on
$\mathcal H_N^-$ they cancel. The negative sector is annihilated,
not mapped by flipped translation.

The translation factor depends on identifying the dual lattice with
the original one. In Ref.~\cite{Aasen2016Topological}, identify the
input spin $h_j$ with $b_j$ and the dual spin
$\hat h_{j+1/2}$ with $a_j$. Its primal-to-dual kernel is then
$\mathcal D_N$, since
\begin{align}
	\sum_j(\hat h_{j-1/2}h_j+h_j\hat h_{j+1/2})
	 & =\sum_j(a_{j-1}b_j+b_ja_j)=q(a,b).
	\notag
\end{align}
Returning by the transpose map gives
$\mathcal D_N^\dagger\mathcal D_N=\Id+\eta$; repeating the same
identified endomorphism gives
$\mathcal D_N^2=(\Id+\eta)T$.
The partial-isometry convention of Ref.~\cite{Seiberg2024Majorana}
has the scale of $U_N$. After exchanging its Pauli $X$ and $Z$,
its displayed intertwining direction is that of $U_N^\dagger$ here:
\begin{align}
	U_N^\dagger X_j        & =Z_jZ_{j+1}U_N^\dagger,
	\notag                                           \\
	U_N^\dagger Z_jZ_{j+1} & =X_{j+1}U_N^\dagger.
	\notag
\end{align}
Accordingly, its translation corresponds to $T^{-1}$ in our
convention, with $(U_N^\dagger)^2=P_+T^{-1}$.
This compares the operator relations and normalization, not an overall
phase in a circuit representation. The noninvertibility here concerns
the periodic bosonic spin chain, not fermionic Majorana translation.
The translation in~(\ref{eq:asymex_kw_fusion}) must be retained when
comparing this lattice relation with a size-independent fusion algebra.

Positive length is essential in these calculations. At $N=1$,
$K_1=\Id+X_0$ and $U_1=P_+$, consistently with the self-bond convention.
The empty virtual word instead has trace $\tr\Id_2=2$, so extending
the raw trace definition to $N=0$ would give $K_0=2$. Its square
would be $4$, whereas the right-hand side of
(\ref{eq:asymex_kw_fusion}) would be $2$.

We now describe the virtual decomposition underlying the square.
For the matrix-unit tensors of Definition~\ref{def:asymex_kw},
expanding the cyclic trace gives
\begin{align}
	\tr(T^{a_0b_0}\cdots T^{a_{N-1}b_{N-1}})
	 & =\prod_j\delta_{a_j,b_{j+1}}=T_{ab},
	\notag                                           \\
	\tr((\eta T)^{a_0b_0}\cdots(\eta T)^{a_{N-1}b_{N-1}})
	 & =\prod_j\delta_{a_j,1-b_{j+1}}=(\eta T)_{ab}.
	\notag
\end{align}
The stacked tensor $B$ generates $K_N^2$: multiplication of periodic
operators sums the intermediate physical index at every site.

\begin{center}
	% Formula: B^{s's}=\sum_mW^{s'm}\otimes W^{ms}, the stacked product
	% of Definition~\ref{def:asymex_kw}.
	% Source: Definition~\ref{def:asymex_kw}.
	% Ink-to-index map: both rows are the duality tensor $W$; the shared
	% vertical pairleg at each site is the summed index $m$; the drawn
	% upper-row label $r$ is the per-site component of the outgoing
	% index $s'$.
	% Contraction set: the vertical pairlegs (index $m$) at every site
	% and the horizontal virtual bonds within each row; the physical
	% legs $(r,s)$, the per-site components of $(s',s)$, remain open.
	% Boundary signature: (virtual-west=0 | virtual-east=0 |
	% physical-up=3 | physical-down=3) on both sides.
	\tenkzkernel
	\begin{tenkz}[rows={op,op}, cols=4, boundary=periodic]
		\tn[skin=mpo, ports={90:physical:$r_1$, 270:physical}]{W} &
		\tn[skin=mpo, ports={90:physical:$r_2$, 270:physical}]{W} &
		\tn[skin=dots, ports={180:virtual, 0:virtual}]{} &
		\tn[skin=mpo, ports={90:physical:$r_N$, 270:physical}]{W} \\
		\tn[skin=mpo, ports={90:physical, 270:physical:$s_1$}]{W} &
		\tn[skin=mpo, ports={90:physical, 270:physical:$s_2$}]{W} &
		\tn[skin=dots, ports={180:virtual, 0:virtual}]{} &
		\tn[skin=mpo, ports={90:physical, 270:physical:$s_N$}]{W}
	\end{tenkz}
	\quad $=$ \quad
	\begin{tenkz}[rows={op}, cols=4, boundary=periodic]
		\tn[skin=mpo, ports={90:physical:$r_1$, 270:physical:$s_1$}]{B} &
		\tn[skin=mpo, ports={90:physical:$r_2$, 270:physical:$s_2$}]{B} &
		\tn[skin=dots, ports={180:virtual, 0:virtual}]{} &
		\tn[skin=mpo, ports={90:physical:$r_N$, 270:physical:$s_N$}]{B}
	\end{tenkz}
\end{center}

The site labels in the diagrams run from $1$ to $N$, a cyclic relabeling
of $I_N$. Normality of a tensor below means that words of some fixed
positive length span its virtual matrix algebra. This differs from
operator normality in~(\ref{eq:asymex_kw_gram}) and does not impose a
normalization on the spectral radius of a transfer map.

\begin{theorem}[Normality of the two translation tensors]
	\label{thm:asymex_kw_normal}
	\lean{KWExample.shiftMPS_isNormal, KWExample.flipShiftMPS_isNormal}
	\leanok
	\uses{def:asymex_kw, def:normal}
	The translation tensor and the flipped translation tensor are normal:
	the four matrices of each are the four matrix units of $\MN{2}$, which
	already span $\MN{2}$ at word length one.
\end{theorem}

\begin{proof}
	\leanok
	Every matrix unit occurs among the four matrices of each tensor, and
	the matrix units span $\MN{2}$. Each tensor is therefore injective
	and hence normal.
\end{proof}

\begin{theorem}[Split compression of the squared duality]
	\label{thm:asymex_kw_compression}
	% Principal declaration: KWExample.kwSquare_isReduction.
	\lean{KWExample.kwSquare_compression, KWExample.kwSquare_isReduction,
		KWExample.kwSquare_left_mul_right_of_ne, KWExample.kwSquare_dim_eq,
		KWExample.kwSquare_remainder_eq_zero}
	\leanok
	\uses{def:asymex_kw, thm:asym_multiblock}
	The stacked tensor of Definition~\ref{def:asymex_kw} has the block
	triangular form of Theorem~\ref{thm:asym_multiblock} for the two
	targets $C_1=2T$ and $C_2=2\eta T$, with no zero slots and $4=2+2$.
	There are compression maps $W_s:\C^4\to\C^2$ and
	$V_s:\C^2\to\C^4$, indexed here by $s=1,2$, such that
	$W_sV_t=\delta_{st}\Id_2$ and $W_sB^wV_s=C_s^w$ for every word $w$.
	The remainder vanishes:
	\begin{align}
		R^{ab} & =B^{ab}-\sum_{s=1}^2V_sC_s^{ab}W_s=0.
		\notag
	\end{align}
	Thus the extension splits.
\end{theorem}

\begin{proof}
	\leanok
	\uses{thm:asymex_explicit_gauge}
	The change of bond coordinates is
	\begin{align}
		G & =\frac12\begin{pmatrix}
			            1 & 0 & 1  & 0  \\
			            0 & 1 & 0  & -1 \\
			            1 & 0 & -1 & 0  \\
			            0 & 1 & 0  & 1
		            \end{pmatrix},
		\label{eq:asymex_kw_gauge}
	\end{align}
	whose inverse has integer entries. For each physical pair $(a,b)$,
	direct multiplication gives
	\begin{align}
		GB^{ab}G^{-1} & =(2T^{ab})\oplus(2(\eta T)^{ab}).
		\notag
	\end{align}
	After clearing the denominator, these are equalities of integer
	matrices. Block diagonality gives the triangularity condition and
	the vanishing remainder. The coordinate projections and inclusions
	give the compression pairs by
	Theorem~\ref{thm:asymex_explicit_gauge}; the dimension count is
	clause (vii) of Theorem~\ref{thm:asym_multiblock}.
\end{proof}

\begin{theorem}[Periodic data of the squared duality]
	\label{thm:asymex_kw_trace}
	\lean{KWExample.kwSquare_trace_evalWord}
	\leanok
	\uses{def:asymex_kw}
	For every nonempty word $w$ over the alphabet of index pairs,
	\begin{align}
		\tr(B^w) & =2^{|w|}(\tr(T^w)+\tr((\eta T)^w)).
		\label{eq:asymex_kw_trace}
	\end{align}
	This is~(\ref{eq:asymex_kw_fusion}) at the level of periodic
	coefficients.
\end{theorem}

\begin{proof}
	\leanok
	\uses{thm:asymex_kw_compression, prop:asym_compression_trace}
	Proposition~\ref{prop:asym_compression_trace} applied to the
	compression of Theorem~\ref{thm:asymex_kw_compression} gives
	$\tr(B^w)=\tr((2T)^w)+\tr((2\eta T)^w)$. Pulling the weight $2$ out
	of each word evaluation contributes the factor $2^{|w|}$.
\end{proof}

\begin{theorem}[The sitewise intertwiners of the split case]
	\label{thm:asymex_kw_intertwiners}
	% Principal declaration: KWExample.kwSquare_mul_right0_eq.
	\lean{KWExample.kwSquareLeft0, KWExample.kwSquareLeft1,
		KWExample.kwSquareRight0, KWExample.kwSquareRight1,
		KWExample.kwSquare_mul_right0_eq, KWExample.kwSquare_mul_right1_eq,
		KWExample.left0_mul_kwSquare_eq, KWExample.left1_mul_kwSquare_eq,
		KWExample.kwSquareLeft0_mul_right0, KWExample.kwSquareLeft1_mul_right1}
	\leanok
	\uses{def:asymex_kw, thm:asymex_kw_compression}
	Write
	\begin{align}
		V_1 & =\begin{pmatrix}
			       1 & 0  \\
			       0 & 1  \\
			       1 & 0  \\
			       0 & -1
		       \end{pmatrix},
		    & V_2
		    & =\begin{pmatrix}
			       1  & 0 \\
			       0  & 1 \\
			       -1 & 0 \\
			       0  & 1
		       \end{pmatrix},
		\notag                        \\
		W_1 & =\frac12\begin{pmatrix}
			              1 & 0 & 1 & 0  \\
			              0 & 1 & 0 & -1
		              \end{pmatrix},
		    & W_2
		    & =\frac12\begin{pmatrix}
			              1 & 0 & -1 & 0 \\
			              0 & 1 & 0  & 1
		              \end{pmatrix}.
		\notag
	\end{align}
	Then $W_1V_1=W_2V_2=\Id_2$, and for every index pair,
	$B^{s's}V_1=V_1(2T)^{s's}$, $B^{s's}V_2=V_2(2\eta T)^{s's}$,
	$W_1B^{s's}=(2T)^{s's}W_1$ and $W_2B^{s's}=(2\eta T)^{s's}W_2$.
\end{theorem}

\begin{proof}
	\leanok
	\uses{thm:asymex_explicit_gauge}
	The matrices $W_1,W_2$ are the two row blocks of $G$, and $V_1,V_2$
	are the corresponding column blocks of $G^{-1}$, as in
	Theorem~\ref{thm:asymex_explicit_gauge}. Each intertwining identity
	is verified by multiplying the explicit matrices and clearing the
	denominator. The two normalizations also follow from clause (iv) of
	Theorem~\ref{thm:asym_multiblock}.
\end{proof}

\begin{center}
	% Formula: the four sitewise identities of
	% Theorem~\ref{thm:asymex_kw_intertwiners}: B^{s's}V_1=V_1(2T)^{s's},
	% B^{s's}V_2=V_2(2\eta T)^{s's}, W_1B^{s's}=(2T)^{s's}W_1, and
	% W_2B^{s's}=(2\eta T)^{s's}W_2.
	% Source: Theorem~\ref{thm:asymex_kw_intertwiners}.
	% Ink-to-index map: the ring is the compression map $V_1$, $W_1$,
	% $V_2$ or $W_2$; the square node is the source $B$ or its target
	% $2T$ or $2\eta T$.
	% Contraction set: the single virtual bond joining the ring to the
	% square node in each pair; the physical pair leg remains open.
	% Boundary signature: (virtual-west=1 | virtual-east=1 |
	% physical-up=1 | physical-down=0) for each pictured equality.
	\tenkzkernel
	\begin{tenkz}[rows={wire}, cols=2, west=open, east=open]
		\tn[label pos=s, ports={90:physical:$s's$}]{B} & \tn[skin=ring]{V_1}
	\end{tenkz}
	\quad $=$ \quad
	\begin{tenkz}[rows={wire}, cols=2, west=open, east=open]
		\tn[skin=ring]{V_1} & \tn[label pos=s, ports={90:physical:$s's$}]{2T}
	\end{tenkz}
	\\
	\begin{tenkz}[rows={wire}, cols=2, west=open, east=open]
		\tn[label pos=s, ports={90:physical:$s's$}]{B} & \tn[skin=ring]{V_2}
	\end{tenkz}
	\quad $=$ \quad
	\begin{tenkz}[rows={wire}, cols=2, west=open, east=open]
		\tn[skin=ring]{V_2} &
		\tn[label pos=s, ports={90:physical:$s's$}]{2\eta T}
	\end{tenkz}
	\\
	\begin{tenkz}[rows={wire}, cols=2, west=open, east=open]
		\tn[skin=ring]{W_1} & \tn[label pos=s, ports={90:physical:$s's$}]{B}
	\end{tenkz}
	\quad $=$ \quad
	\begin{tenkz}[rows={wire}, cols=2, west=open, east=open]
		\tn[label pos=s, ports={90:physical:$s's$}]{2T} & \tn[skin=ring]{W_1}
	\end{tenkz}
	\\
	\begin{tenkz}[rows={wire}, cols=2, west=open, east=open]
		\tn[skin=ring]{W_2} & \tn[label pos=s, ports={90:physical:$s's$}]{B}
	\end{tenkz}
	\quad $=$ \quad
	\begin{tenkz}[rows={wire}, cols=2, west=open, east=open]
		\tn[label pos=s, ports={90:physical:$s's$}]{2\eta T} &
		\tn[skin=ring]{W_2}
	\end{tenkz}
\end{center}

\section{Kramers--Wannier duality acting on ground states}%
\label{sec:asymex_kw_action}

The finite-chain actions in~(\ref{eq:asymex_kw_state_actions}) have
different virtual decompositions. The product-state action splits into
the two GHZ sectors with a local weight $2$. The GHZ action has two
unweighted copies of the same product-state target, together with a
nonzero remainder. Product and GHZ tensors are elementary MPS examples
in Ref.~\cite{Cirac2021Matrix}; here their explicit matrices allow the
asymmetric compression and its sitewise intertwiners to be compared
directly.

\begin{definition}[The two action tensors]
	\label{def:asymex_kw_action}
	\lean{KWExample.plusTensor, KWExample.ghz0, KWExample.ghz1,
		KWExample.kwPlus, KWExample.plusTarget, KWExample.ghz,
		KWExample.kwGHZ, KWExample.kwGHZTarget}
	\leanok
	\uses{def:asymex_kw, def:mps_tensor}
	The product-state tensor $A_+$ has bond dimension one and
	$A_+^0=A_+^1=(1)$. Its periodic vector is $p_N$.
	The two ferromagnetic-sector tensors have
	$C_0^i=(\delta_{i,0})$ and $C_1^i=(\delta_{i,1})$.
	Their direct sum is the GHZ tensor
	$A_{\mathrm{GHZ}}^0=\diag(1,0)$,
	$A_{\mathrm{GHZ}}^1=\diag(0,1)$, whose periodic vector is $g_N$.

	For either input tensor $A$, the action tensor is
	$(W\cdot A)^{s'}=\sum_sW^{s's}\otimes A^s$.
	For $A=A_+$ its matrices are
	\begin{align}
		(W\cdot A)^0
		 & =\begin{pmatrix}1&1\\1&1\end{pmatrix},
		 &
		(W\cdot A)^1
		 & =\begin{pmatrix}1&-1\\-1&1\end{pmatrix},
		\notag
	\end{align}
	with targets $2C_0$ and $2C_1$. For $A=A_{\mathrm{GHZ}}$ they are
	\begin{align}
		(W\cdot A)^0
		 & =\begin{pmatrix}
			    1 & 0 & 1 & 0 \\
			    0 & 0 & 0 & 0 \\
			    0 & 0 & 0 & 0 \\
			    0 & 1 & 0 & 1
		    \end{pmatrix},
		 &
		(W\cdot A)^1
		 & =\begin{pmatrix}
			    1 & 0  & -1 & 0 \\
			    0 & 0  & 0  & 0 \\
			    0 & 0  & 0  & 0 \\
			    0 & -1 & 0  & 1
		    \end{pmatrix},
		\notag
	\end{align}
	with two unweighted copies of $A_+$ as targets.
\end{definition}

\begin{center}
	% Formula: (W\cdot A)^{s'}=\sum_sW^{s's}\otimes A^s, of
	% Definition~\ref{def:asymex_kw_action}, evaluated at the
	% product-state tensor $A_+$ (left) and at the Greenberger--Horne--
	% Zeilinger tensor $A_{\mathrm{GHZ}}$ (right).
	% Source: Definition~\ref{def:asymex_kw_action}.
	% Ink-to-index map: the upper row is the duality tensor $W$, whose
	% drawn physical legs $r$ are the per-site components of the
	% outgoing index $s'$; the lower row is the acted-on state tensor;
	% the shared vertical pairleg at each site is the summed index $s$.
	% Contraction set: the vertical pairlegs (index $s$) at every site
	% and the horizontal virtual bonds within each row, both closed by
	% the periodic trace; only the upper physical legs of $W$ remain
	% open.
	% Boundary signature: (virtual-west=0 | virtual-east=0 |
	% physical-up=3 | physical-down=0) for both pictures.
	\tenkzkernel
	\begin{tenkz}[rows={op,ket}, cols=4, boundary=periodic]
		\tn[skin=mpo, ports={90:physical:$r_1$}]{W} &
		\tn[skin=mpo, ports={90:physical:$r_2$}]{W} &
		\tn[skin=dots, ports={180:virtual, 0:virtual}]{} &
		\tn[skin=mpo, ports={90:physical:$r_N$}]{W} \\
		\tn[label pos=s, ports={180:virtual, 0:virtual}]{A_+} &
		\tn[label pos=s, ports={180:virtual, 0:virtual}]{A_+} &
		\tn[skin=dots, ports={180:virtual, 0:virtual}]{} &
		\tn[label pos=s, ports={180:virtual, 0:virtual}]{A_+}
	\end{tenkz}
	\qquad
	\begin{tenkz}[rows={op,ket}, cols=4, boundary=periodic]
		\tn[skin=mpo, ports={90:physical:$r_1$}]{W} &
		\tn[skin=mpo, ports={90:physical:$r_2$}]{W} &
		\tn[skin=dots, ports={180:virtual, 0:virtual}]{} &
		\tn[skin=mpo, ports={90:physical:$r_N$}]{W} \\
		\tn[label pos=s, ports={180:virtual, 0:virtual}]{A_{\mathrm{GHZ}}} &
		\tn[label pos=s, ports={180:virtual, 0:virtual}]{A_{\mathrm{GHZ}}} &
		\tn[skin=dots, ports={180:virtual, 0:virtual}]{} &
		\tn[label pos=s, ports={180:virtual, 0:virtual}]{A_{\mathrm{GHZ}}}
	\end{tenkz}
\end{center}

For the product-state action, the matrices are $2Q_+$ and $2Q_-$,
where
\begin{align}
	Q_\pm & =\frac12\begin{pmatrix}1&\pm1\\\pm1&1\end{pmatrix}.
	\notag
\end{align}
They satisfy $Q_\pm^2=Q_\pm$, $Q_+Q_-=Q_-Q_+=0$ and $\tr Q_\pm=1$.
Hence a nonempty word of length $N$ has trace $2^N$ when all letters
agree and trace zero otherwise. This recovers $K_Np_N=2^Ng_N$ and
identifies the factor $2$ as a local tensor weight. By contrast, the
factor $2$ in $K_Ng_N=2p_N$ is the multiplicity of two unweighted
targets. Replacing $p_N$ or $g_N$ by a normalized vector changes the
global scalar as in Lemma~\ref{lem:asymex_kw_states}; it does not
change the weights of these fixed tensors.

In the following two results, $B$ denotes the relevant action tensor,
rather than the stacked square tensor of the preceding section.

\begin{theorem}[The duality on the product state splits]
	\label{thm:asymex_kw_plus}
	% Principal declaration: KWExample.plus_isReduction.
	\lean{KWExample.plusCompression, KWExample.plus_isReduction,
		KWExample.plus_dim_eq, KWExample.plus_remainder_eq_zero,
		KWExample.kwPlus_trace_evalWord_eq_sum}
	\leanok
	\uses{def:asymex_kw_action, thm:asym_multiblock}
	The first action tensor of Definition~\ref{def:asymex_kw_action} has
	the block triangular form of Theorem~\ref{thm:asym_multiblock} for
	the weighted targets $\widehat C_s=2C_s$, $s=0,1$, with no zero
	slots and $2=1+1$. Its compression pairs satisfy $W_sV_s=\Id_1$
	and $W_sB^wV_s=\widehat C_s^w$ for every word $w$.
	The remainder vanishes identically, and for every nonempty word,
	\begin{align}
		\tr(B^w) & =\tr((2C_0)^w)+\tr((2C_1)^w).
		\notag
	\end{align}
\end{theorem}

\begin{proof}
	\leanok
	\uses{thm:asymex_explicit_gauge, prop:asym_compression_trace}
	Let $F=\begin{pmatrix}1&1\\1&-1\end{pmatrix}$, so
	$F^{-1}=F/2$. Direct multiplication gives
	\begin{align}
		FB^0F^{-1} & =\diag(2,0),
		\notag                    \\
		FB^1F^{-1} & =\diag(0,2).
		\notag
	\end{align}
	Diagonality gives the triangularity condition, identifies the
	matched blocks, and makes the remainder zero. The coordinate
	compression pairs satisfy the stated word identities by
	Theorem~\ref{thm:asymex_explicit_gauge}. The trace identity follows
	from Proposition~\ref{prop:asym_compression_trace}.
\end{proof}

\begin{theorem}[The duality on the ferromagnetic state does not split]
	\label{thm:asymex_kw_ghz}
	% Principal declaration: KWExample.kwGHZ_isReduction.
	\lean{KWExample.kwGHZCompression, KWExample.kwGHZ_isReduction,
		KWExample.kwGHZ_dim_eq, KWExample.kwGHZ_evalWord_remainder_eq_zero,
		KWExample.kwGHZRemainder_eq, KWExample.kwGHZ_remainder_mul_remainder,
		KWExample.kwGHZ_trace_evalWord_eq_two}
	\leanok
	\uses{def:asymex_kw_action, thm:asym_multiblock}
	The second action tensor of Definition~\ref{def:asymex_kw_action}
	has the block triangular form of Theorem~\ref{thm:asym_multiblock}
	for two unweighted copies $C_s=A_+$ of the product-state tensor,
	with $z=2$ zero slots and $4=1+1+2$. The compression pair of each
	copy satisfies $W_sV_s=\Id_1$ and $W_sB^wV_s=C_s^w$ for every word.
	For every nonempty word, $\tr(B^w)=2$.
	Every remainder word of length at least four vanishes; in fact,
	\begin{align}
		R^0 & =\begin{pmatrix}
			       0 & 0 & 1 & 0 \\
			       0 & 0 & 0 & 0 \\
			       0 & 0 & 0 & 0 \\
			       0 & 1 & 0 & 0
		       \end{pmatrix},
		    & R^1
		    & =\begin{pmatrix}
			       0 & 0  & -1 & 0 \\
			       0 & 0  & 0  & 0 \\
			       0 & 0  & 0  & 0 \\
			       0 & -1 & 0  & 0
		       \end{pmatrix},
		\notag
	\end{align}
	and $R^iR^j=0$ for all letters $i,j$.
\end{theorem}

\begin{proof}
	\leanok
	\uses{thm:asymex_explicit_gauge, prop:asym_compression_trace}
	The two copies occupy the first and fourth bond coordinates, and
	the zero slots occupy the second and third. Ordering the coordinates
	as $(0,3,1,2)$ makes both letters upper triangular, with diagonal
	$(1,1,0,0)$. Thus the matched blocks are both $A_+$, and the
	coordinate compression pairs give the asserted reductions.
	Every word evaluation of $A_+$ is $1$, so
	Proposition~\ref{prop:asym_compression_trace} gives
	$\tr(B^w)=1+1=2$ for nonempty $w$.

	Subtracting the matched diagonal contributions in the original
	coordinates gives the displayed remainder matrices. Their only
	nonzero entries lie in positions $(0,2)$ and $(3,1)$, so direct
	multiplication gives $R^iR^j=0$. The general bound in clause (vi)
	of Theorem~\ref{thm:asym_multiblock} also gives vanishing at lengths
	at least $|S|+z=4$, where $S$ is the set of two matched slots.
	The obstruction to a sitewise splitting is established in
	Theorem~\ref{thm:asymex_kw_ghz_intertwiner}.
\end{proof}

A word reduction need not provide a pair of sitewise intertwiners.
Here the periodic coefficients have merged into two copies of the same
product state, but the virtual action still has a one-sided obstruction.

\begin{theorem}[One-sided sitewise intertwiners for the merged target]
	\label{thm:asymex_kw_ghz_intertwiner}
	\lean{KWExample.kwGHZ_left_intertwiner_eq_zero,
		KWExample.kwGHZ_right_intertwiner,
		KWExample.kwGHZ_right_intertwiner_ne_zero}
	\leanok
	\uses{def:asymex_kw_action}
	Let $B=W\cdot A_{\mathrm{GHZ}}$. If a row vector $u\in\C^4$
	satisfies $uB^i=u$ for both letters $i$, then $u=0$.
	Thus no nonzero sitewise left intertwiner exists for the
	product-state target. A nonzero sitewise right intertwiner does
	exist: the first coordinate vector satisfies $B^ie_0=e_0$ for both
	letters.
\end{theorem}

\begin{proof}
	\leanok
	Both target matrices are the scalar $1$. In coordinates
	$u=(u_0,u_1,u_2,u_3)$, the left-intertwiner equations are
	\begin{align}
		uB^0 & =(u_0,u_3,u_0,u_3)=u,
		\notag                         \\
		uB^1 & =(u_0,-u_3,-u_0,u_3)=u.
		\notag
	\end{align}
	The second coordinates give $u_3=u_1=0$, and the third give
	$u_0=u_2=0$. For the right intertwiner, the first column of each
	matrix is $e_0$, which is nonzero.
\end{proof}

This obstruction rules out a sitewise split copy of the product-state
target: such a splitting would require a nonzero left intertwiner.
It does not obstruct the word reductions or the periodic trace formula
in Theorem~\ref{thm:asymex_kw_ghz}. Those statements retain the
finite-chain action $K_Ng_N=2p_N$ while distinguishing it from a split
virtual direct sum.
